% ════════════════════════════════════════════════════════════════════════════
%  Database Architecture Demonstration Packet
%    §4.1 Database Identity Model
%    §4.2 Input Dimensionality and Data Typing
%    §4.3 Element, Isotope, and Material Shard Layout
%    §4.4 Compound and Precursor Registry Design
%    §4.5 Seed and Hash Provenance
%    §4.6 Integrated Mixed-Resolution Demonstration Packet
%  Reports 1–10 of the Day 40C Report Suite.
%  Printer-ready LaTeX.  Compile with pdflatex (twice for refs).
% ════════════════════════════════════════════════════════════════════════════
\documentclass[11pt, a4paper, oneside]{article}

% ── Geometry ────────────────────────────────────────────────────────────────
\usepackage[
  top=1.0in, bottom=1.0in, left=1.15in, right=1.15in,
  headheight=14pt
]{geometry}

% ── Fonts / encoding ───────────────────────────────────────────────────────
\usepackage[T1]{fontenc}
\usepackage{lmodern}
\usepackage{microtype}

% ── Math ───────────────────────────────────────────────────────────────────
\usepackage{amsmath, amssymb, amsthm, mathtools}
\usepackage{bm}

% ── Tables ─────────────────────────────────────────────────────────────────
\usepackage{booktabs}
\usepackage{array}
\usepackage{tabularx}
\usepackage{longtable}

% ── Lists ──────────────────────────────────────────────────────────────────
\usepackage{enumitem}
\setlist{nosep, leftmargin=1.5em}

% ── Code ───────────────────────────────────────────────────────────────────
\usepackage{listings}
\usepackage{xcolor}

\lstset{
  basicstyle=\ttfamily\small,
  keywordstyle=\color{blue},
  commentstyle=\color{gray},
  stringstyle=\color{red!70!black},
  breaklines=true,
  columns=fullflexible,
  keepspaces=true,
  frame=single,
  backgroundcolor=\color{gray!8},
  xleftmargin=0.5em,
  xrightmargin=0.5em,
  language=C++
}

\lstdefinelanguage{JSON}{
  basicstyle=\ttfamily\small,
  string=[s]{"}{"},
  stringstyle=\color{red!70!black},
  comment=[l]{//},
  commentstyle=\color{gray},
  morecomment=[s]{/*}{*/},
  literate=
    *{:}{{{\color{blue}:}}}{1}
     {,}{{{\color{blue},}}}{1}
     {\{}{{{\color{blue}\{}}}{1}
     {\}}{{{\color{blue}\}}}}{1}
     {[}{{{\color{blue}[}}}{1}
     {]}{{{\color{blue}]}}}{1}
}

% ── Headers / footers ─────────────────────────────────────────────────────
\usepackage{fancyhdr}
\pagestyle{fancy}
\fancyhf{}
\fancyhead[L]{\small\textit{VSEPR-SIM Database Architecture}}
\fancyhead[R]{\small\textit{Demonstration Packet — Day 40C}}
\fancyfoot[C]{\thepage}
\renewcommand{\headrulewidth}{0.4pt}

% ── Section formatting ────────────────────────────────────────────────────
\usepackage{titlesec}
\titleformat{\section}{\Large\bfseries}{§\thesection}{0.8em}{}
\titleformat{\subsection}{\large\bfseries}{§\thesubsection}{0.6em}{}
\titleformat{\subsubsection}{\normalsize\bfseries\itshape}{§\thesubsubsection}{0.5em}{}

% ── Custom commands ───────────────────────────────────────────────────────
\newcommand{\field}[1]{\texttt{#1}}
\newcommand{\dtype}[1]{\texttt{#1}}
\newcommand{\recdot}{\,\cdot\,}

% ════════════════════════════════════════════════════════════════════════════
\begin{document}

\begin{center}
{\LARGE\bfseries Database Architecture Demonstration Packet}\\[0.6em]
{\large VSEPR-SIM Day 40C Report Suite — Reports 1\,--\,10}\\[0.4em]
{\normalsize\itshape Typed Scientific Registry for Atomistic,
  Material, Compound, and Precursor Records}\\[1.2em]
\rule{0.7\textwidth}{0.5pt}
\end{center}
\vspace{1em}

\tableofcontents
\newpage

% ════════════════════════════════════════════════════════════════════════════
\section{Database Architecture Demonstration Reports}
\label{sec:db_arch}
% ════════════════════════════════════════════════════════════════════════════

The database layer is treated as a typed scientific registry rather than a flat
storage table.  Each input is modeled as
\begin{equation}
\label{eq:generic_record}
X_i = (I_i,\; T_i,\; D_i,\; P_i,\; S_i,\; H_i),
\end{equation}
where $I_i$ is the identity block, $T_i$ is the type block, $D_i$ is the
dimension descriptor, $P_i$ is the scientific payload, $S_i$ is the
deterministic seed block, and $H_i$ is the hash and provenance block.

The architecture is designed to host:
\begin{itemize}
  \item elemental records,
  \item compound and material records,
  \item precursor staging objects,
  \item mixed-resolution simulation inputs,
  \item deterministic reconstruction using seed and hash provenance.
\end{itemize}

The flow is:
\[
\text{Input Layer}
\;\longrightarrow\;
\text{Canonical Registry}
\;\longrightarrow\;
\text{Projection / Resolution Layer}
\;\longrightarrow\;
\text{Validation / Hash / Report Layer}.
\]


% ════════════════════════════════════════════════════════════════════════════
\section{Report 1: Input Dimension Demonstration}
\label{sec:report1}
% ════════════════════════════════════════════════════════════════════════════

\subsection{Goal}

This report demonstrates how simulation and database inputs are described in
terms of dimensionality, resolution level, physical scale, storage type, and
projection target.  The intent is to ensure that every input object has a known
shape before entering either the atomistic engine, the bead engine, or the
formation engine.

\subsection{General Input Form}

For a generic scientific record $X_i$, the six blocks of
Eq.~\eqref{eq:generic_record} expand as follows:

\begin{center}
\begin{tabular}{@{} cl @{}}
\toprule
\textbf{Block} & \textbf{Description}\\
\midrule
$I_i$ & Identity block (name, symbol, canonical ID)\\
$T_i$ & Type block (element, compound, material, precursor, \ldots)\\
$D_i$ & Dimension block (see below)\\
$P_i$ & Payload block (physical / chemical / nuclear data)\\
$S_i$ & Seed / deterministic initialization block\\
$H_i$ & Hash / provenance block\\
\bottomrule
\end{tabular}
\end{center}

\subsection{Dimension Block}
\label{sec:dim_block}

A dimension descriptor is written as
\begin{equation}
\label{eq:dim_block}
D_i = \bigl(n_{\text{state}},\; n_{\text{geom}},\; n_{\text{chem}},\;
             n_{\text{nuc}},\; n_{\text{aux}}\bigr),
\end{equation}
allowing a single registry to host element cores, isotope packets, atomistic
objects, bead objects, materials, compounds, and staged precursors.

\begin{center}
\begin{tabular}{@{} cl @{}}
\toprule
\textbf{Component} & \textbf{Meaning}\\
\midrule
$n_{\text{state}}$ & Number of runtime state variables\\
$n_{\text{geom}}$  & Spatial / geometric dimensions\\
$n_{\text{chem}}$  & Chemical descriptor count\\
$n_{\text{nuc}}$   & Nuclear descriptor count\\
$n_{\text{aux}}$   & Metadata / diagnostics / tags\\
\bottomrule
\end{tabular}
\end{center}

\subsection{Input Dimension Table}
\label{sec:dim_table}

\begin{center}
\small
\begin{tabular}{@{} l l c c c l @{}}
\toprule
\textbf{Input Class} & \textbf{Repr.} & $n_{\text{geom}}$
  & $n_{\text{chem}}$ & $n_{\text{nuc}}$ & \textbf{Typical Payload}\\
\midrule
Element core      & scalar record  & 0   & 4--12  & 0--4   & name, symbol, $Z$\\
Isotope record    & scalar record  & 0   & 2--8   & 4--12  & decay, half-life\\
Atomistic object  & vector state   & 3   & 6--20  & 0--8   & $\bm{r},\bm{v},q,\sigma_{\text{chem}},\sigma_{\text{nuc}}$\\
Bead object       & reduced state  & 3   & 4--16  & 0--6   & $R,V,Q,\Theta,\beta,\chi$\\
Compound record   & graph / set    & 0--3 & 8--64  & 0--12  & stoichiometry, class\\
Material preset   & fielded record & 0--3 & 10--48 & 0--8   & density, phase, modulus\\
Precursor object  & staged object  & 0--3 & 6--32  & 0--8   & reactive flags, seed geometry\\
\bottomrule
\end{tabular}
\end{center}

\subsection{Demonstration Inputs}

\paragraph{Example A: Hydrogen element core.}
\[
X_H^{\text{core}} = \bigl(Z{=}1,\;\text{name}{=}\text{Hydrogen},\;
  \text{symbol}{=}\text{H},\;\text{class}{=}\text{nonmetal}\bigr).
\]

\paragraph{Example B: Plutonium atomistic object.}
\[
X_{\text{Pu}}^{(a)} = (\bm{r},\;\bm{v},\;q,\;\sigma_{\text{chem}},\;\sigma_{\text{nuc}}),
\]
with
\[
\sigma_{\text{nuc}} = (\lambda,\;E_{\text{dec}},\;N_{\text{emit}},\;\tau,\;\text{mode}).
\]

\paragraph{Example C: Cyclopentadienyl precursor bead.}
\[
X_{\text{Cp}}^{(b)} = (R,\;V,\;Q,\;\Theta,\;\beta,\;\chi,\;\eta).
\]

This demonstrates that the database accepts multiple dimensional classes without
pretending they are all the same object.


% ════════════════════════════════════════════════════════════════════════════
\section{Report 2: Data Type Demonstration}
\label{sec:report2}
% ════════════════════════════════════════════════════════════════════════════

\subsection{Goal}

This report demonstrates the internal data typing rules used by the
architecture.  Every field is assigned a \emph{semantic type} and a
\emph{storage type}.

\subsection{Typed Record Form}

A field descriptor is:
\begin{equation}
\label{eq:field_desc}
F_{ij} = \bigl(\text{name},\;\text{semantic\_type},\;\text{storage\_type},\;
               \text{units},\;\text{validity}\bigr).
\end{equation}

\subsection{Semantic Type Classes}

\begin{center}
\begin{tabular}{@{} l l l @{}}
\toprule
\textbf{Semantic Type} & \textbf{Meaning} & \textbf{Example}\\
\midrule
\dtype{identity}    & Immutable object identifier      & $Z{=}94$, Pu\\
\dtype{categorical} & Discrete class label             & actinide, ligand, precursor\\
\dtype{scalar}      & Single real/integer value         & mass, radius, charge\\
\dtype{vector}      & Ordered numeric list              & position, velocity\\
\dtype{tensor}      & Matrix-like or directional data   & inertia, polarizability\\
\dtype{set}         & Unordered discrete collection     & oxidation states\\
\dtype{graph}       & Connectivity relation             & bonding map\\
\dtype{hash}        & Content signature                 & SHA-256 / custom 562\\
\dtype{pathref}     & Path to canonical shard           & \texttt{/elements/094\_Pu/core.json}\\
\bottomrule
\end{tabular}
\end{center}

\subsection{Storage Type Classes}

\begin{center}
\begin{tabular}{@{} l l @{}}
\toprule
\textbf{Storage Type} & \textbf{Use}\\
\midrule
\dtype{u32}       & Compact indices, element IDs\\
\dtype{u64}       & Extended object IDs\\
\dtype{f32}       & Fast approximate simulation values\\
\dtype{f64}       & Scientific numeric values\\
\dtype{str}       & Names, symbols, tags\\
\dtype{bool}      & Flags, response switches\\
\dtype{vec<T>}    & Ordered arrays\\
\dtype{set<T>}    & Non-duplicate category collections\\
\dtype{map<K,V>}  & Keyed lookup data\\
\dtype{hash256}   & Record integrity (SHA-256)\\
\dtype{hash562}   & Deep provenance digest (custom)\\
\bottomrule
\end{tabular}
\end{center}

\subsection{Demo Schema}

\begin{center}
\small
\begin{tabular}{@{} l l l l l @{}}
\toprule
\textbf{Field} & \textbf{Semantic} & \textbf{Storage} & \textbf{Units} & \textbf{Validity}\\
\midrule
\field{atomic\_number}     & identity  & \dtype{u32}        & ---   & immutable\\
\field{half\_life\_s}      & scalar    & \dtype{f64}        & s     & isotope-only\\
\field{oxidation\_states}  & set       & \dtype{vec<i32>}   & ---   & context-sensitive\\
\field{bead\_orientation}  & tensor    & \dtype{vec<f64>}   & rad   & bead projection only\\
\bottomrule
\end{tabular}
\end{center}


% ════════════════════════════════════════════════════════════════════════════
\section{Report 3: Material Demonstration Report}
\label{sec:report3}
% ════════════════════════════════════════════════════════════════════════════

\subsection{Goal}

This report demonstrates how engineering and chemistry materials are stored as
typed overlays rather than collapsed into element records.

A material is not an element.

\subsection{Material Record Form}

\begin{equation}
\label{eq:material_record}
M_k = \bigl(I_k,\; C_k,\; \rho_k,\; \phi_k,\; \Pi_k,\; \Gamma_k,\; H_k\bigr),
\end{equation}
where:
\begin{center}
\begin{tabular}{@{} cl @{}}
\toprule
\textbf{Block} & \textbf{Description}\\
\midrule
$I_k$        & Material identity\\
$C_k$        & Composition (stoichiometric or mass-fraction)\\
$\rho_k$     & Density / bulk density\\
$\phi_k$     & Phase state\\
$\Pi_k$      & Engineering property vector\\
$\Gamma_k$   & Hazard / environment descriptor\\
$H_k$        & Provenance hash block\\
\bottomrule
\end{tabular}
\end{center}

\subsection{Example Materials}

\paragraph{A.\;Copper.}
\[
M_{\text{Cu}} = \bigl(\text{Copper},\;\{\text{Cu}{:}1.0\},\;
  \rho,\;\text{solid},\;\Pi,\;\Gamma,\;H\bigr).
\]

\paragraph{B.\;Steel preset (A36).}
\[
M_{\text{Steel-A36}} = \bigl(\text{A36},\;\{\text{Fe},\text{C},\text{Mn},\ldots\},\;
  \rho,\;\text{solid},\;E,\sigma_y,\nu,\ldots\bigr).
\]

\paragraph{C.\;Plutonium oxide-bearing residue.}
\[
M_{\text{PuResidue}} = \bigl(\text{mixed residue},\;
  \{\text{Pu},\text{O},\text{C},\text{H},\text{trace}\},\;
  \rho,\;\phi,\;\Pi,\;\Gamma,\;H\bigr).
\]

\subsection{Demo Material Table}

\begin{center}
\small
\begin{tabular}{@{} l l l l @{}}
\toprule
\textbf{Material} & \textbf{Class} & \textbf{Phase} & \textbf{Main Use}\\
\midrule
Copper                & elemental material         & solid     & conductor, engineering\\
Iron                  & elemental material         & solid     & base metal\\
Uranium dioxide       & compound material          & solid     & nuclear fuel model\\
Thorium oxalate       & precursor / compound       & solid     & slurry / precipitation / staging\\
Hexene                & organic precursor          & liquid    & bead-scale precursor\\
Cyclopentadienyl prec.& organometallic ligand      & molecular & hybrid architecture tests\\
\bottomrule
\end{tabular}
\end{center}


% ════════════════════════════════════════════════════════════════════════════
\section{Report 4: Compound Demonstration Report}
\label{sec:report4}
% ════════════════════════════════════════════════════════════════════════════

\subsection{Goal}

This report demonstrates how compounds are encoded as distinct database objects
with stoichiometric, structural, and resolution-aware metadata.

\subsection{Compound Record Form}

\begin{equation}
\label{eq:compound_record}
C_j = \bigl(S_j,\; G_j,\; B_j,\; P_j,\; R_j,\; H_j\bigr),
\end{equation}
where:
\begin{center}
\begin{tabular}{@{} cl @{}}
\toprule
\textbf{Block} & \textbf{Description}\\
\midrule
$S_j$ & Stoichiometry\\
$G_j$ & Geometry / graph / arrangement hints\\
$B_j$ & Bond or interaction classification\\
$P_j$ & Phase / process metadata\\
$R_j$ & Resolution projections\\
$H_j$ & Compound hash block\\
\bottomrule
\end{tabular}
\end{center}

\subsection{Demonstration Compounds}

\paragraph{A.\;Water.}
\[
C_{\text{H}_2\text{O}} =
  \bigl(\{\text{H}{:}2,\text{O}{:}1\},\;\text{bent},\;\text{polar covalent},\;
        \text{molecular},\;R,\;H\bigr).
\]

\paragraph{B.\;Thorium oxalate.}
\[
C_{\text{Th-ox}} =
  \bigl(\{\text{Th},\text{C},\text{O},\text{H}\},\;\text{coordination / precipitation},\;
        \text{ionic-covalent hybrid},\;\text{precursor solid},\;R,\;H\bigr).
\]

\paragraph{C.\;Uranium dioxide.}
\[
C_{\text{UO}_2} =
  \bigl(\{U{:}1,O{:}2\},\;\text{crystal lattice},\;\text{ionic / lattice},\;
        \text{ceramic solid},\;R,\;H\bigr).
\]

\paragraph{D.\;Hexene.}
\[
C_{\text{C}_6\text{H}_{12}} =
  \bigl(\{C{:}6,H{:}12\},\;\text{chain / isomer family},\;\text{covalent organic},\;
        \text{molecular liquid},\;R,\;H\bigr).
\]

\subsection{Compound Demo Payload}

\begin{lstlisting}[language=JSON, caption={Hexene compound registry entry}]
{
  "compound_id":   "cmp_hexene_001",
  "name":          "hexene",
  "formula":       "C6H12",
  "class":         "organic_precursor",
  "stoichiometry": {"C": 6, "H": 12},
  "structure_family": "alkene",
  "projection_modes": ["atomistic", "bead"],
  "default_seed_class": "chain_seed_v1",
  "hash256": "....",
  "hash562": "...."
}
\end{lstlisting}


% ════════════════════════════════════════════════════════════════════════════
\section{Report 5: Precursor Demonstration Report}
\label{sec:report5}
% ════════════════════════════════════════════════════════════════════════════

\subsection{Goal}

This report demonstrates how precursors are treated as staged formation objects
rather than fully committed final structures.  A precursor is not just a
compound.  It is a reaction-positioned entity.

\subsection{Precursor Record Form}

\begin{equation}
\label{eq:precursor_record}
P_m = \bigl(I_m,\; \Psi_m,\; \Omega_m,\; \Lambda_m,\; T_m,\; H_m\bigr),
\end{equation}
where:
\begin{center}
\begin{tabular}{@{} cl @{}}
\toprule
\textbf{Block} & \textbf{Description}\\
\midrule
$I_m$        & Precursor identity\\
$\Psi_m$     & Chemical state\\
$\Omega_m$   & Environment readiness\\
$\Lambda_m$  & Transformation / reaction flags\\
$T_m$        & Target formation class\\
$H_m$        & Provenance block\\
\bottomrule
\end{tabular}
\end{center}

\subsection{Example Precursors}

\begin{center}
\small
\begin{tabular}{@{} l l l @{}}
\toprule
\textbf{Precursor} & \textbf{Class} & \textbf{Target Formation}\\
\midrule
Hexene bead prec.               & organic chain precursor      & bead-tested hydrocarbon family\\
Cyclopentadienyl ligand prec.   & organometallic precursor     & mixed bead/atomistic coordination\\
Thorium oxalate slurry prec.    & precipitation precursor      & oxide or residue evolution\\
Dissolved nitrate--metal complex & process precursor            & sludge / salt / oxide outcome\\
Pipe-scale dissolved solids     & process-field precursor      & materials + deposition model\\
\bottomrule
\end{tabular}
\end{center}

\subsection{Hybrid Precursor Demo}

A mixed-resolution precursor packet:
\begin{equation}
\label{eq:hybrid_prec}
P_{\text{hybrid}} = \bigl(X_{\text{Pu}}^{(a)},\;
  X_{\text{Cp}}^{(b)},\; X_{\text{hexene}}^{(b)},\; E,\; T,\; H\bigr),
\end{equation}
where $X_{\text{Pu}}^{(a)}$ is atomistic plutonium,
$X_{\text{Cp}}^{(b)}$ is bead-scale cyclopentadienyl precursor,
$X_{\text{hexene}}^{(b)}$ is bead-scale hexene precursor,
$E$ is the environment state, and $T$ is the transformation intent.


% ════════════════════════════════════════════════════════════════════════════
\section{Report 6: Seed and Hash Demonstration Report}
\label{sec:report6}
% ════════════════════════════════════════════════════════════════════════════

\subsection{Goal}

This report demonstrates deterministic reconstruction and provenance control.
Every major input carries both a seed for reproducible generation and a hash for
identity and integrity verification.

\subsection{Seed Block}

\begin{equation}
\label{eq:seed_block}
S_i = \bigl(\text{seed}_{32},\;\text{seed}_{64},\;
            \text{seed}_{\text{geom}},\;\text{seed}_{\text{chem}}\bigr),
\end{equation}
where seeds may control:
\begin{itemize}
  \item randomized cloud generation,
  \item bead placement,
  \item lattice jitter,
  \item initial velocity assignment,
  \item synthetic precursor arrangement,
  \item test case regeneration.
\end{itemize}

\subsection{Hash Block}

\begin{equation}
\label{eq:hash_block}
H_i = \bigl(H_{32}^{\text{tag}},\; H_{256}^{\text{content}},\;
            H_{562}^{\text{provenance}}\bigr).
\end{equation}

Recommended roles:
\begin{center}
\begin{tabular}{@{} c l @{}}
\toprule
\textbf{Hash} & \textbf{Role}\\
\midrule
$H_{32}^{\text{tag}}$         & Quick human-friendly tag\\
$H_{256}^{\text{content}}$    & Exact record content fingerprint (SHA-256)\\
$H_{562}^{\text{provenance}}$ & Deep lineage signature (custom)\\
\bottomrule
\end{tabular}
\end{center}

\subsection{Demonstration Table}

\begin{center}
\small
\begin{tabular}{@{} l c c c @{}}
\toprule
\textbf{Object} & \textbf{Seed} & \textbf{Content Hash} & \textbf{Provenance Hash}\\
\midrule
H element core                & fixed canonical & yes & yes\\
Pu isotope packet             & optional        & yes & yes\\
Hexene bead scene             & required        & yes & yes\\
Thorium oxalate test cloud    & required        & yes & yes\\
Hybrid Pu--Cp--Hexene system  & required        & yes & yes\\
\bottomrule
\end{tabular}
\end{center}

\subsection{Demo Entry}

\begin{lstlisting}[language={}, caption={Seed and hash block for hybrid scene}]
Object:       hybrid_precursor_scene_004
Seed32:       40128
Seed64:       40128_094_006_005
GeometrySeed: geom_cp_hexene_pu_v2
ChemicalSeed: chem_hybrid_transition_v1
Hash256:      98C1...A442
Hash562:      HYB-PU-CP-HEX-2026-TRACE-....
\end{lstlisting}


% ════════════════════════════════════════════════════════════════════════════
\section{Report 7: Integrated Example Report}
\label{sec:report7}
% ════════════════════════════════════════════════════════════════════════════

\subsection{Goal}

This report shows one full input packet combining element identity, material
typing, compound registration, precursor staging, and seed and hash
traceability.

\subsection{Example: Pu + Cyclopentadienyl + Hexene Hybrid Input}

\paragraph{Input packet.}
\begin{equation}
\label{eq:demo_packet}
I_{\text{demo}} = \bigl(E_{\text{Pu}},\; C_{\text{Cp-like}},\;
  C_{\text{hexene}},\; P_{\text{hybrid}},\; S,\; H\bigr).
\end{equation}

\paragraph{Expanded form.}
\begin{align*}
E_{\text{Pu}}        &= (I,\;B,\;N,\;R,\;C,\;X,\;H),\\
C_{\text{Cp-like}}   &= (S,\;G,\;B,\;P,\;R,\;H),\\
P_{\text{hybrid}}    &= (I,\;\Psi,\;\Omega,\;\Lambda,\;T,\;H).
\end{align*}

\subsection{Demo Manifest}

\begin{lstlisting}[language={}, caption={Integrated demonstration manifest}]
Demo Name: Pu_Cp_Hexene_Hybrid_001
Mode: mixed_resolution
Primary Atomistic Core: Pu
Supporting Precursors: cyclopentadienyl_like, hexene
Environment Class: dry_condensed_transition
Dimensions:
  state_dim = 41
  geom_dim  = 3
  chem_dim  = 19
  nuc_dim   = 5
  aux_dim   = 12

Data Types:
  core_id         : u32
  runtime_id      : u64
  positions       : f64[3]
  charges         : f64
  oxidation_set   : i32[]
  decay_constants : f64[]
  flags           : bool[]
  labels          : str[]
  hash256         : bytes32
  hash562         : custom_trace

Materials:
  plutonium_core
  organic_precursor_family
  organometallic_transition_packet

Compounds:
  hexene
  cyclopentadienyl_like

Precursors:
  Cp_precursor_stage2
  Hexene_bead_stage1

Seed:
  seed32    = 41092
  seed64    = 41092_2026_40C_001
  geom_seed = ring_chain_offset_v1
  chem_seed = hybrid_decay_response_v1

Hash:
  hash256 = ...
  hash562 = ...
\end{lstlisting}


% ════════════════════════════════════════════════════════════════════════════
\section{Report 8: Recommended Chapter Placement}
\label{sec:report8}
% ════════════════════════════════════════════════════════════════════════════

If this material becomes part of a formal textbook or architecture document, the
recommended placement is:

\begin{center}
\begin{tabular}{@{} l l @{}}
\toprule
\textbf{Section} & \textbf{Topic}\\
\midrule
§4.1 & Database identity model\\
§4.2 & Input dimensionality and data typing\\
§4.3 & Element, isotope, and material shard layout\\
§4.4 & Compound and precursor registry design\\
§4.5 & Seed and hash provenance\\
§4.6 & Integrated mixed-resolution demonstration packet\\
\bottomrule
\end{tabular}
\end{center}


% ════════════════════════════════════════════════════════════════════════════
\section{Report 9: Clean \LaTeX{} Subsection Starter}
\label{sec:report9}
% ════════════════════════════════════════════════════════════════════════════

The following verbatim block is the canonical \LaTeX{} subsection starter text
for embedding in other documents:

\begin{lstlisting}[language=TeX, caption={Embeddable \LaTeX{} subsection starter}]
\section{Database Architecture Demonstration Reports}

The database layer is treated as a typed scientific registry
rather than a flat storage table. Each input is modeled as
\[
X_i = (I_i, T_i, D_i, P_i, S_i, H_i),
\]
where \(I_i\) is the identity block, \(T_i\) is the type block,
\(D_i\) is the dimension descriptor, \(P_i\) is the scientific
payload, \(S_i\) is the deterministic seed block, and \(H_i\) is
the hash and provenance block.

\subsection{Input Dimension Demonstration}
A dimension descriptor is written as
\[
D_i = (n_{\text{state}}, n_{\text{geom}}, n_{\text{chem}},
       n_{\text{nuc}}, n_{\text{aux}}),
\]
allowing a single registry to host element cores, isotope packets,
atomistic objects, bead objects, materials, compounds, and staged
precursors.

\subsection{Seed and Hash Provenance}
To guarantee reproducibility and traceability, each major object
carries both a deterministic seed block and a multi-scale hash block:
\[
S_i = (\text{seed}_{32}, \text{seed}_{64},
       \text{seed}_{geom}, \text{seed}_{chem}),
\]
\[
H_i = (H_{32}^{tag}, H_{256}^{content}, H_{562}^{provenance}).
\]
This allows exact regeneration of synthetic scenes, validation of
imported records, and lineage tracing across atomistic, bead, and
macro-scale projections.
\end{lstlisting}


% ════════════════════════════════════════════════════════════════════════════
\section{Report 10: Best Next Expansions}
\label{sec:report10}
% ════════════════════════════════════════════════════════════════════════════

The next useful reports building on this architecture would be:

\begin{enumerate}
  \item \textbf{Element demo packet} — H, Fe, U, Pu with full
    $X_i = (I_i, T_i, D_i, P_i, S_i, H_i)$ records, including
    nuclear descriptor expansions for actinides.

  \item \textbf{Materials demo packet} — Copper, Steel~A36, UO\textsubscript{2},
    Thorium oxalate with $M_k$ records and engineering property vectors.

  \item \textbf{Precursor demo packet} — Hexene, cyclopentadienyl,
    nitrate--metal slurries, dissolved crystal solids, with thermal
    and strength initial passes for tensile hardness and brittleness
    if those channels are validated.

  \item \textbf{Seed reproducibility report} — Demonstrate that the same
    seed produces the same scene.  Show 3D Qt6 rendering with at least
    two physical representations per atom.  Print additional math steps
    for the reconstruction chain.

  \item \textbf{Hash mismatch report} — Demonstrate detection of tampered
    overlays, comparing $H_{256}^{\text{content}}$ before and after
    modification and flagging provenance violations.
\end{enumerate}


% ════════════════════════════════════════════════════════════════════════════
\section*{Appendix: C++ Header Mapping}
\label{sec:appendix_headers}
% ════════════════════════════════════════════════════════════════════════════

The formal definitions in this document are implemented as C++ header-only
structures under \texttt{coarse\_grain/database/}:

\begin{center}
\begin{tabular}{@{} l l @{}}
\toprule
\textbf{Header} & \textbf{Implements}\\
\midrule
\texttt{input\_record.hpp}     & $X_i$, $D_i$, \texttt{SemanticType}, \texttt{StorageType}\\
\texttt{material\_record.hpp}  & $M_k$\\
\texttt{compound\_record.hpp}  & $C_j$\\
\texttt{precursor\_record.hpp} & $P_m$\\
\texttt{seed\_hash.hpp}        & $S_i$, $H_i$\\
\texttt{demo\_entries.hpp}     & Canonical demonstration objects\\
\bottomrule
\end{tabular}
\end{center}

% ════════════════════════════════════════════════════════════════════════════
\end{document}
